\naslov{Manjkajoče vrednosti}

Včasih v testnih ali učnih primerih manjka kakšen podatek.
Če manjka $y$, primer samo izbrišemo, ker nam ne pomaga.
To lahko naredimo tudi za manjkajoče $X_i$, čemur pravimo \pojem{complete-case
  analysis}.
Tu lahko pride do problemov; eno je predsodek v podatkih, drugo pa, da ima tako
učna množica lahko premalo primerov.
Podobna ideja je \pojem{available-case analysis}, kjer namesto primerov brišemo
spremenljivke (stolpce), kjer podatki manjkajo.

\vprasanje{Kakšne so težave z brisanjem neznanih podatkov?}

Če podatkov ne brišemo, delamo \pojem{imputacijo}.
Numerično spremenljivko lahko nadomestimo s povprečjem vseh primerov v
podatkovni množici, diskretno pa z gostiščnico, ali pa neznane podatke postavimo
v razred zase.
Dobra ideja je, da dodamo še novo indikatorsko spremenljivko, ki pove, ali je
bila originalna spremenljivka prisotna, ali pa je imputirana.
Druga možnost, ki je baje \enquote{relativno korektna,} je, da neznano vrednost
nadomestimo z naključno vrednostjo iz tega stolpca.
To je dejanje iz obupa; v praksi je bolje, da uporabiš znanje o problemu.
Korektno pa je, ker tako ne prinašaš predsodka, ker se porazdelitev ne spremeni.
Še zadnja osnovna možnost je nadomeščanje s povprečno vrednostjo najbližjih
sosedov (po znanih spremenljivkah).

\vprasanje{Opiši osnovne tehnike imputacije.}

Bolj napredno je nadomeščanje z napovednimi modeli.
To je iterativni postopek, kjer v prvem koraku napovemo manjkajoče vrednosti z
eno od enostavnih metod, v vseh nadaljnjih korakih pa izberemo eno od
spremenljivk (v nekem vrstnem redu), naučimo napovedni model na vseh $X_i$ brez
te, in napovemo vrednost te spremenljivke.
Tu \poudarjenne{} uporabljamo podatkov iz $Y$, ker bi potem bile generirane
vrednosti pristranske, kar nam efektivno uniči testne primere.

\vprasanje{Opiši imputacijo z napovednimi modeli.}

Izjema tu so modeli, ki sprejemajo tudi neznane vrednosti, recimo odločitvena
drevesa.
V algoritmu TDIDT primere z neznano vrednostjo $X_i$ enakomerno porazdelimo med
vse naslednike; te primere v računih obravnavamo, kot da bi bili na pol v eni in
na pol v drugi množici.

\vprasanje{Kako prilagodiš TDIDT tako, da lahko dela tudi z neznanimi
  vrednostmi?}

% LocalWords:  imputacijo complete-case analysis available-case imputirana
% LocalWords:  imputacije generirane TDIDT
